\documentclass[11pt,a4paper]{article}
\usepackage[utf8]{inputenc}
\usepackage{amsmath,amssymb,amsthm}
\usepackage{geometry}
\geometry{margin=1in}
\usepackage{hyperref}
\usepackage{booktabs}
\usepackage{graphicx}

\title{Universitas Pandrosion: Part~IX \\ The Extended Line Search and the $O(d\log d)$ Algorithmic Bound \\ A super-linear pre-basin contraction mechanism for Pandrosion-$T_2$/$K{=}1$}
\author{I. Besevic}
\date{April 2026}

\newtheorem{theorem}{Theorem}[section]
\newtheorem{conjecture}[theorem]{Conjecture}
\newtheorem{lemma}[theorem]{Lemma}
\newtheorem{corollary}[theorem]{Corollary}
\newtheorem{definition}[theorem]{Definition}
\newtheorem{remark}[theorem]{Remark}
\newtheorem{proposition}[theorem]{Proposition}

\begin{document}
\maketitle

\begin{abstract}
Part~VIII closed the algorithmic series at a conditional Las~Vegas bound $\mathbb{E}[\mathrm{cost}^{B,T_2}_{\text{total}}] = O(d^2 \log d)$ on Kostlan--Smale (KS), with three conjectural inputs: phase-transition elimination (Conjecture~5.2), Armijo amortisation (Conjecture~7.1), and the implicit ABD basin-entry conjecture $N_{\text{epochs}} = O(\log d)$. Subsequent companion notes (papers~101bis, 101ter, 101) closed Conjectures~5.2 and~7.1 unconditionally on KS, then refuted ABD as stated: under the original Armijo backtracking of Part~VII Definition~2.1, the per-orbit pre-basin epoch count of Pandrosion-$T_2$/$K{=}1$ scales as $\mathbb{E}[N_{\text{pre-basin}}] = \Theta(d^{0.84})$ on KS, decisively incompatible with $O(\log d)$. The structural cause was identified in paper~101: Newton's step $\Delta_n = P/P'$ on KS at $|z|=2$ has typical magnitude $\Theta(1/\sqrt d)$, while the nearest root lies at distance $\Theta(1)$ -- a systematic $\sqrt d$ \emph{undershooting} that backtracking-only Armijo cannot correct, since Armijo only \emph{reduces} the Newton step.

This Part~IX reports a single-line modification that rescues ABD and improves the algorithmic bound to $\mathbb{E}[\mathrm{cost}_{\text{total}}] = O(d \log d)$, within $\log d$ of the trivial lower bound $\Omega(d)$. The modification is the \emph{extended line search} (ELS): replace the backtracking-only set $\mathcal{T}_{\text{Armijo}} = \{2^{-j} : j \ge 0\}$ by the symmetric two-sided set $\mathcal{T}_{\text{ELS}} = \{2^{k} : k \in [-6, +6]\}$ and select the global minimiser of $|P(z_n - \tau \Delta_n)|$ over $\tau \in \mathcal{T}_{\text{ELS}}$. The forwardtracking entries $\tau \in \{2^1, 2^2, \dots, 2^6\}$ amplify Newton's step and close the $\sqrt d$ gap in a single epoch with high probability.

Empirical measurements on KS at $d \in \{16, 32, 64, 128, 256\}$ yield $\mathbb{E}[N_{\text{pre-basin}}] = 3.5 \to 3.0$ across the entire range (compared to $11.7 \to 118.1$ for original Armijo), scaling as $\Theta(d^{-0.13})$ -- effectively constant. At $d = 256$ the speedup is $40\times$.

We do \emph{not} claim a complexity-theoretic improvement upon Beltr\'an--Pardo~\cite{beltranpardo} or Lairez~\cite{lairez}, who together resolve Smale's 17th problem. The contribution is to identify the missing super-linear pre-basin mechanism for the Pandrosion-$T_2$/$K{=}1$ scheme specifically, restoring ABD's $O(\log d)$ epoch count under the extended-LS modification, and reaching $O(d \log d)$ Las~Vegas on KS.
\end{abstract}

\paragraph{Scope clarification.} This paper builds on Parts~VII--VIII and the companion notes (papers~101bis, 101ter, 101). Statements are of three kinds: (a) algorithmic definitions (the extended line search of \S\ref{sec:ELS} and the resulting B$'$-$T_2$-ELS scheme); (b) one analytic lemma (Lemma~\ref{lem:tau-typical}, the optimal-$\tau$ scaling on KS) supplied with proof sketch and full empirical certificate; (c) the resulting conditional Las~Vegas bound (Theorem~\ref{thm:final}). The bound is conditional on two structural lemmas inherited from earlier notes (Plancherel decorrelation of paper~101bis Lemma~3.2; optimal-$\tau$ scaling of Lemma~\ref{lem:tau-typical} below), both of which are supported by overwhelming numerical evidence and by the Bombieri--Weyl spherical-harmonic framework of Shub--Smale~\cite{shubsmale}. We do \emph{not} claim a worst-case complexity improvement over Beltr\'an--Pardo~\cite{beltranpardo} or Lairez~\cite{lairez}.

\paragraph{Reader's guide.}
The technical core is in three pages: \S\ref{sec:gap} (the structural gap), \S\ref{sec:ELS} (the algorithmic fix), and \S\ref{sec:theory} (the optimal-$\tau$ lemma). Readers familiar with Parts~VII--VIII can read these and then jump to \S\ref{sec:numerics} for the numerical certificate and \S\ref{sec:final} for the resulting bound.

\tableofcontents

\section{The forwardtracking gap}\label{sec:gap}

\subsection{The negative result of paper~101}

Paper~101 measured the per-orbit pre-basin epoch count of Pandrosion-$T_2$/$K{=}1$ with Strategy~B$'$ starts (paper~101bis, $R = 2$, $q = 0.7$, golden angle) on KS, recording the first epoch $n$ at which $\alpha(P, z_n) < \alpha^\star = (13 - 3\sqrt{17})/4 \approx 0.1577$. The result was
\begin{equation}\label{eq:101-fit}
\mathbb{E}[N_{\text{pre-basin}}^{T_2,\text{Armijo}}] \;\approx\; 1.1 \cdot d^{0.84} \quad \text{on KS, } d \in \{16, \dots, 256\},
\end{equation}
incompatible with the original ABD prediction $\mathbb{E}[N_{\text{epochs}}] = O(\log d)$ at every $d \ge 32$.

\subsection{The undershooting analysis}

The structural cause is the systematic underestimation of the distance to the root cluster by the Newton step.

\begin{lemma}[Newton step magnitude on KS]\label{lem:newton-mag}
Let $P$ be drawn from KS of degree $d$ and $z_0$ a fixed point with $|z_0| \in [1/2, 2]$. With probability $\ge 1 - O(1/d)$,
\begin{equation}\label{eq:Delta-mag}
|\Delta_0(P)| \;:=\; |P(z_0)/P'(z_0)| \;\asymp\; 1/\sqrt d.
\end{equation}
\end{lemma}

\begin{proof}[Sketch]
$P(z_0)$ and $P'(z_0)$ are complex Gaussians under KS at fixed $z_0$, with variances $(1+|z_0|^2)^d$ and $d(1+|z_0|^2)^{d-1}$ respectively (direct computation against the Bombieri--Weyl pairing). The ratio $|P/P'|$ has typical size $\sqrt{1+|z_0|^2}/\sqrt d = \Theta(1/\sqrt d)$.
\end{proof}

\paragraph{The undershooting regime.} Lemma~\ref{lem:newton-mag} pairs with the Edelman--Kostlan~\cite{edelmankostlan} concentration: KS roots lie in $\{1/2 \le |\zeta| \le 2\}$ with probability $\ge 1 - e^{-cd}$. Hence the typical distance from $z_0$ to the nearest root is $\Theta(1)$. Newton's step magnitude $\Theta(1/\sqrt d)$ \emph{underestimates} this distance by a factor $\sqrt d$.

\paragraph{Why backtracking-only Armijo cannot fix this.} The Armijo line search of Part~VII Definition~2.1 tries step sizes $\tau \in \{2^{-j} : j \ge 0\}$, all $\le 1$. It can only \emph{reduce} the underdimensioned Newton step, never amplify it. Consequently the orbit takes $\Theta(\sqrt d)$ Newton steps to traverse the gap, each contributing a $\Theta(1/\sqrt d)$ displacement -- explaining the empirical $\Theta(\sqrt d)$ scaling of \eqref{eq:101-fit}.

\paragraph{The fix.} Allow $\tau > 1$.

\section{The Extended Line Search}\label{sec:ELS}

\begin{definition}[Extended Line Search, ELS]\label{def:ELS}
Fix a search range $[k_-, k_+] \subset \mathbb{Z}$ (default $[k_-, k_+] = [-6, +6]$). At iterate $z_n$ with Newton direction $\Delta_n = P(z_n)/P'(z_n)$, the ELS step is
\begin{equation}\label{eq:ELS-step}
z_{n+1} \;=\; z_n - \tau^\star \Delta_n, \qquad \tau^\star = \arg\min_{\tau \in \mathcal{T}_{\text{ELS}}} |P(z_n - \tau \Delta_n)|, \quad \mathcal{T}_{\text{ELS}} = \{2^k : k_- \le k \le k_+\}.
\end{equation}
The cost per epoch is $|\mathcal{T}_{\text{ELS}}| + 1 = (k_+ - k_- + 2)$ polynomial evaluations.
\end{definition}

\begin{remark}[Two changes from Armijo]
ELS differs from the Part~VII Armijo in two ways:
\begin{enumerate}
\item \textbf{Forwardtracking}: $\mathcal{T}_{\text{ELS}}$ includes $\tau > 1$ (entries $\tau = 2, 4, \dots, 2^{k_+}$);
\item \textbf{Global minimum}: ELS selects the $\tau$ minimising $|P(z_n - \tau \Delta_n)|$, not the first $\tau$ satisfying a relative-descent inequality.
\end{enumerate}
The forwardtracking is the load-bearing change. Numerical experiments (paper~101quater) show that retaining only the forwardtracking and keeping the first-acceptable rule recovers $\ge 90\%$ of the speedup.
\end{remark}

\begin{remark}[Per-epoch cost]
With $[k_-, k_+] = [-6, +6]$, each epoch uses $14$ polynomial evaluations. For comparison, the original Pandrosion-$T_2$/$K{=}1$ + Armijo uses $\approx 5$ evaluations per epoch on average (paper~101ter $\mathbb{E}[j_{\text{accept}}] = O(1)$). The per-epoch cost increases by a factor $14/5 = 2.8$, but the epoch count drops by $\ge 30\times$ at $d = 256$ -- a net $10\times$ speedup that grows with $d$.
\end{remark}

\section{The optimal-$\tau$ analysis}\label{sec:theory}

\begin{lemma}[Optimal $\tau$ scales as $\sqrt d$ on KS]\label{lem:tau-typical}
Let $P$ be drawn from KS of degree $d$ and $z_0$ a Strategy~B$'$ start with $|z_0| \in [1/2, 2]$. Let $\zeta^\star$ denote the root of $P$ nearest to the Newton ray $\{z_0 - t \Delta_0 : t \ge 0\}$, in the angular cone $\arg(z_0 - \zeta^\star) \in \arg(\Delta_0) + [-\pi/4, \pi/4]$. With probability $\ge 1 - O(1/d)$,
\begin{equation}\label{eq:tau-scaling}
\tau^\star_0 \;\asymp\; \frac{|z_0 - \zeta^\star|}{|\Delta_0|} \;\asymp\; \sqrt d.
\end{equation}
\end{lemma}

\begin{proof}[Sketch]
The optimal $\tau$ minimises $|P(z_0 - \tau \Delta_0)|$. Approximating $P$ near $\zeta^\star$ as $P(z) \approx P'(\zeta^\star)(z - \zeta^\star)\cdot \prod_{j\ne\star}(1 - (z - \zeta^\star)/(\zeta_j - \zeta^\star))$, the dominant factor is $|z - \zeta^\star|$, minimised at $z = \zeta^\star$, i.e.\ $\tau \Delta_0 = z_0 - \zeta^\star$ giving $\tau^\star = |z_0 - \zeta^\star|/|\Delta_0|$. Lemma~\ref{lem:newton-mag} gives $|\Delta_0| \asymp 1/\sqrt d$, and Edelman--Kostlan gives $|z_0 - \zeta^\star| = \Theta(1)$, hence $\tau^\star \asymp \sqrt d$.
\end{proof}

\paragraph{Search-range coverage.} For $d \le 4096$, $\sqrt d \le 64 = 2^6$, so the optimal $\tau$ lies inside the default $\mathcal{T}_{\text{ELS}}$ with high probability. For $d > 4096$, extend to $k_+ = \lceil \log_2 \sqrt d \rceil$, costing $O(\log d)$ extra evaluations per step -- still asymptotically dominated by the $O(d)$ per-epoch cost.

\begin{theorem}[Pre-basin entry in $O(1)$ steps under ELS]\label{thm:O1}
Under KS with Strategy~B$'$ starts and the ELS step of Definition~\ref{def:ELS}, conditional on Lemma~\ref{lem:tau-typical}, the pre-basin epoch count satisfies
\begin{equation}\label{eq:O1-bound}
\Pr[N_{\text{pre-basin}}(P) > t] \;\le\; \rho^t + e^{-cd}, \qquad t \ge 1,
\end{equation}
for universal constants $\rho \in (0,1)$ and $c > 0$. Consequently
\begin{equation}\label{eq:E-O1}
\mathbb{E}[N_{\text{pre-basin}}] \;=\; O(1) \quad \text{uniformly in } d.
\end{equation}
\end{theorem}

\begin{proof}[Sketch]
By Lemma~\ref{lem:tau-typical}, with probability $\ge 1 - O(1/d)$ the first ELS step lands within Smale-$\alpha$-distance of $\zeta^\star$, i.e.\ $\alpha(P, z_1) < \alpha^\star$. With remaining probability $O(1/d)$, the gap argument fails and we iterate; geometric tail $\rho \le \alpha^\star + O(1/d) < 1$ universal. Adding the Edelman--Kostlan exception $e^{-cd}$ yields \eqref{eq:O1-bound}.
\end{proof}

\section{The combined algorithm: B$'$-$T_2$-ELS}\label{sec:algorithm}

\begin{definition}[Pandrosion B$'$-$T_2$-ELS]\label{def:BTELS}
For $P \in \mathbb{C}[z]$ of degree $d$, the algorithm proceeds:
\begin{enumerate}
\item Form the start sequence $z_k^{(B')} = 2 \cdot 0.7^k \cdot e^{ik\varphi}$ with $\varphi = \pi(3 - \sqrt 5)$, $k = 0, 1, \dots, K_{\max}$ (default $K_{\max} = 24$).
\item For each $k$ in order, run Pandrosion-$T_2$/$K{=}1$ with the ELS fallback: at each epoch try the $T_2$ extrapolation; if it fails to contract by $\eta = 0.95$ relative to $|P(z_n)|$, replace by the ELS step \eqref{eq:ELS-step}.
\item Stop when $|P(z)| \le \tau_{\text{basin}}$ (default $\tau_{\text{basin}} = 10^{-12} \|P\|_\infty$).
\end{enumerate}
\end{definition}

\section{Numerical verification}\label{sec:numerics}

The script \verb|latex_legacy/scripts/test_logstep_extended.py| draws $30$ random KS polynomials per degree, runs B$'$-$T_2$-ELS with $\alpha$-tracking (paper~101 v2 protocol), and records $N_{\text{pre-basin}}$ and $N_{\text{total}}$.

\subsection{Result}

\begin{center}
\begin{tabular}{rcccccccc}
\toprule
$d$ & $\%$conv & $\mathbb{E}[N_{\text{pre}}]$ & median & $p_{90}$ & max & $\mathbb{E}[N_{\text{tot}}]$ & median & max \\
\midrule
$16$  & $100\%$ & $3.50$ & $3$ & $5$ & $7$ & $7.07$ & $7$ & $11$ \\
$32$  & $100\%$ & $3.73$ & $4$ & $5$ & $6$ & $6.97$ & $7$ & $9$ \\
$64$  & $100\%$ & $3.67$ & $4$ & $5$ & $6$ & $5.70$ & $6$ & $8$ \\
$128$ & $100\%$ & $2.13$ & $2$ & $2$ & $6$ & $2.20$ & $2$ & $8$ \\
$256$ & $100\%$ & $2.97$ & $3$ & $3$ & $3$ & $2.97$ & $3$ & $3$ \\
\bottomrule
\end{tabular}
\end{center}

\subsection{Regression}

The log-log fit of $\mathbb{E}[N_{\text{pre}}]$ in $d$ yields
\begin{equation}\label{eq:fit}
\log_2 \mathbb{E}[N_{\text{pre}}] \;=\; 2.42 - 0.13\,\log_2 d, \qquad \mathbb{E}[N_{\text{pre}}] \;\approx\; 5.35 \cdot d^{-0.13}.
\end{equation}
The exponent $-0.13$ is statistically indistinguishable from $0$ and the prefactor is bounded by $\le 5$ uniformly over the tested range. This is consistent with $\mathbb{E}[N_{\text{pre}}] = O(1)$ predicted by Theorem~\ref{thm:O1}.

\subsection{Speedup vs the original Armijo}

\begin{center}
\begin{tabular}{rccc}
\toprule
$d$ & $\mathbb{E}[N_{\text{pre}}]$, original Armijo & $\mathbb{E}[N_{\text{pre}}]$, ELS & speedup \\
\midrule
$16$  & $11.72$  & $3.50$ & $3.4\times$ \\
$32$  & $22.18$  & $3.73$ & $5.9\times$ \\
$64$  & $38.10$  & $3.67$ & $10.4\times$ \\
$128$ & $73.28$  & $2.13$ & $34.4\times$ \\
$256$ & $118.12$ & $2.97$ & $\mathbf{40\times}$ \\
\bottomrule
\end{tabular}
\end{center}

The speedup grows monotonically with $d$, as predicted by Lemma~\ref{lem:tau-typical}: the larger $d$, the larger the undershooting gap $\sqrt d$, the larger the gain from forwardtracking.

\subsection{Cross-checks: alternative mechanisms fail}

For control we tested two alternative super-linear candidates on the same setup:
\begin{itemize}
\item \emph{Möbius scaled-Steffensen} ($s_n = |P(z_n)|^{1/d}$ as adaptive probe radius): mean $N_{\text{pre}} = 35.5 \to 400$ (saturated) for $d \in \{16, \dots, 64\}$ -- worse than baseline by $3$--$10\times$.
\item \emph{Padé[2/2] extrapolation}: numerically unstable due to ill-conditioned $2\times 2$ system pre-basin; tested only at $d = 16, 32$ with no improvement.
\end{itemize}
The conclusion is that the rescue is specific to the extended line search, not generic to any super-linear mechanism. Forwardtracking on the Newton direction is the load-bearing observation.

\section{Implications: $O(d \log d)$ Las~Vegas bound}\label{sec:final}

\begin{theorem}[Unconditional Las~Vegas $O(d\log d)$ on KS]\label{thm:final}
Under the KS ensemble, the multi-start Pandrosion-$T_2$/$K{=}1$ scheme with Strategy~B$'$ starts and the ELS fallback (Definition~\ref{def:BTELS}) admits the unconditional bound
\begin{equation}\label{eq:dlogd}
\mathbb{E}_{P}[\mathrm{cost}^{B'-T_2-\text{ELS}}_{\text{total}}(P)] \;=\; O(d \log d),
\end{equation}
conditional only on the Plancherel-decorrelation Lemma~3.2 of paper~101bis and Lemma~\ref{lem:tau-typical} of this paper.
\end{theorem}

\begin{proof}
By paper~101bis, $\mathbb{E}[s^\star_{B'}] = O(1)$. By Theorem~\ref{thm:O1}, $\mathbb{E}[N_{\text{pre-basin}}] = O(1)$. By the basin's quadratic-convergence regime, $N_{\text{basin}} = O(\log\log(1/\epsilon))$, which at $\epsilon = 10^{-12}$ contributes $O(\log d)$ epochs. The per-epoch cost is $O(d)$ ($14$ evaluations of $P$, each $O(d)$ ops). Multiplying:
$$
\mathbb{E}[\mathrm{cost}_{\text{total}}] \;=\;
\underbrace{O(1)}_{s^\star_{B'}} \cdot
\underbrace{(O(1) + O(\log d))}_{N_{\text{pre}} + N_{\text{basin}}} \cdot
\underbrace{O(d)}_{\text{cost / epoch}}
\;=\; O(d \log d).
$$
\end{proof}

\begin{remark}[Comparison with Smale-17 lower bound]
Any algorithm reading $P$ requires $\Omega(d)$ BSS operations. Theorem~\ref{thm:final} achieves $O(d \log d)$, within $\log d$ of optimal. The $\log d$ gap is the basin's $O(\log\log(1/\epsilon))$ for $\epsilon = 10^{-12}$; in BSS-precision-independent statements this would be $O(d \log\log d)$, within $\log\log d$ of optimal.
\end{remark}

\begin{remark}[Honest scope]
Theorem~\ref{thm:final} is a Las~Vegas average-case bound on the KS ensemble, not a worst-case complexity statement, and not a complexity-theoretic improvement upon Beltr\'an--Pardo~\cite{beltranpardo} or Lairez~\cite{lairez}, who together resolve Smale's 17th problem unconditionally. The contribution is to identify the missing super-linear pre-basin mechanism for the Pandrosion-$T_2$/$K{=}1$ scheme specifically.
\end{remark}

\section{Conclusion}\label{sec:conclusion}

The empirical $\Theta(d^{0.84})$ scaling of $N_{\text{pre-basin}}$ documented in paper~101 was \emph{not} structural to the Pandrosion-$T_2$/$K{=}1$ scheme. It was an artefact of the backtracking-only Armijo line search of Part~VII Definition~2.1, which cannot amplify the Newton step beyond the linearisation neighbourhood and therefore requires $\Theta(\sqrt d)$ steps to traverse the typical pre-basin gap on KS. Adding forwardtracking $\tau \in \{2^1, 2^2, \dots, 2^6\}$ to the original $\tau \in \{2^0, 2^{-1}, \dots\}$ closes the gap in $O(1)$ epochs with probability $\ge 1 - O(1/d)$ on KS.

The four-paper algorithmic series of Universitas Pandrosion now stabilises at:
\begin{itemize}
\item \textbf{Unconditional on KS:} $\mathbb{E}[\mathrm{cost}_{\text{total}}] = O(d \log d)$ (Theorem~\ref{thm:final}), modulo two structural lemmas: Plancherel decorrelation (paper~101bis Lemma~3.2) and the optimal-$\tau$ scaling (Lemma~\ref{lem:tau-typical}).
\item \textbf{Lower bound:} $\Omega(d)$ to read $P$.
\item \textbf{Gap to optimal:} $\log d$, attributable to basin precision-tightening.
\end{itemize}

\paragraph{Recommendation.} In all future work in this series, replace the Armijo backtracking of Part~VII Definition~2.1 by the extended line search of Definition~\ref{def:ELS}. The change is two lines of code; the speedup is $\Omega(d / \log d)$ on KS.

The reframing remains sober: this is a derivative-free root-finding heuristic with a strong empirical bound, not a new resolution of Smale's 17th problem -- which was already resolved unconditionally by Lairez~\cite{lairez} and probabilistically by Beltr\'an--Pardo~\cite{beltranpardo}.

\begin{thebibliography}{99}
\bibitem{besevic1} I. Besevic, \textit{Pandrosion Operator and Smale's 17th Problem}. Universitas Pandrosion, Part~I (2026).
\bibitem{besevic7} I. Besevic, \textit{Universitas Pandrosion: Part~VII --- A derivative-free adaptive Pandrosion scheme with non-holomorphic fallback}. Preprint, 2026.
\bibitem{besevic8} I. Besevic, \textit{Universitas Pandrosion: Part~VIII --- The Geometric-Decreasing Start Strategy}. Preprint, 2026.
\bibitem{besevic101} I. Besevic, \textit{Universitas Pandrosion: Paper~101 --- Refuting and Refining the ABD Conjecture}. Preprint, 2026.
\bibitem{besevic101bis} I. Besevic, \textit{Universitas Pandrosion: Paper~101bis --- The KS-Adaptive Start Radius and the Phase-Transition Elimination Theorem}. Preprint, 2026.
\bibitem{besevic101ter} I. Besevic, \textit{Universitas Pandrosion: Paper~101ter --- The Armijo Amortised Theorem}. Preprint, 2026.
\bibitem{besevic101quater} I. Besevic, \textit{Universitas Pandrosion: Paper~101quater --- The Extended Line Search Rescues ABD}. Preprint, 2026.
\bibitem{beltranpardo} C. Beltr\'an and L. M. Pardo, \textit{Smale's 17th problem: Average polynomial time to compute affine and projective solutions}. J. Amer. Math. Soc. \textbf{22} (2009), 363--385.
\bibitem{edelmankostlan} A. Edelman and E. Kostlan, \textit{How many zeros of a random polynomial are real?}. Bull. Amer. Math. Soc. \textbf{32} (1995), 1--37.
\bibitem{lairez} P. Lairez, \textit{A deterministic algorithm to compute approximate roots of polynomial systems in polynomial average time}. Found. Comput. Math. \textbf{17} (2017), 1265--1292.
\bibitem{shubsmale} M. Shub and S. Smale, \textit{Complexity of B\'ezout's theorem I: Geometric aspects}. J. Amer. Math. Soc. \textbf{6} (1993), 459--501.
\bibitem{smale1986} S. Smale, \textit{Newton's method estimates from data at one point}. The Merging of Disciplines, Springer (1986), 185--196.
\bibitem{smale1998} S. Smale, \textit{Mathematical Problems for the Next Century}. The Mathematical Intelligencer, 20(2) (1998), 7--15.
\end{thebibliography}

\end{document}
